\documentclass[aspectratio=169,11pt]{beamer}

\usepackage[indonesian]{babel}
\usepackage{amsmath,booktabs,array,tabularx,tikz,listings,xcolor}
\usetikzlibrary{arrows.meta,calc,positioning,decorations.pathreplacing}
\usetheme{metropolis}
\metroset{progressbar=frametitle,block=fill}
\setbeamertemplate{title page}{%
  {\usebeamerfont{title}\inserttitle\par}
  {\usebeamerfont{subtitle}\insertsubtitle\par}
  \usebeamertemplate{title separator}\medskip
  {\usebeamerfont{author}\insertauthor\par}
  {\usebeamerfont{date}\insertdate\par}}

\setbeamerfont{title}{size=\Large}
\setbeamerfont{subtitle}{size=\normalsize}
\usebeamercolor{normal text}\colorlet{codeink}{normal text.fg}
\usebeamercolor{alerted text}\colorlet{codeaccent}{alerted text.fg!60!codeink}

\lstset{
  language={[LaTeX]TeX},
  basicstyle=\ttfamily\scriptsize\color{codeink},
  texcsstyle=*\bfseries\color{codeaccent},
  columns=fullflexible,
  keepspaces=true,
  showstringspaces=false,
  breaklines=true,
  backgroundcolor=\color{codeink!6},
  frame=l,
  rulecolor=\color{codeaccent},
  xleftmargin=5pt,
  aboveskip=4pt,
  belowskip=4pt
}
\lstdefinestyle{plaincode}{basicstyle=\ttfamily\small\color{codeink},breaklines=true}
\newcommand{\tugas}[1]{\par\medskip\alert{\textbf{Coba sendiri}}\quad #1}
\newcommand{\dimensi}[1]{\texttt{\detokenize{#1}}}

\title{Tata Letak Dokumen \LaTeX}
\subtitle{Hari 5 --- Margin, padding, spacing, dan layout}
\author{Pelatihan LaTeX}
\date{21 November 2026}

\begin{document}

\begin{frame}[plain]
  \titlepage
\end{frame}

\begin{frame}{Tujuan belajar Hari 5}
  Pada akhir sesi, peserta dapat:
  \begin{enumerate}
    \item membaca model geometris halaman dan membedakan margin dari area teks;
    \item mengatur ukuran halaman, margin, header, footer, dan ruang jilid;
    \item membedakan \emph{padding}, \emph{gap}, indentasi, dan jarak antarbaris;
    \item menyusun teks dengan panjang relatif seperti \dimensi{\textwidth} dan \dimensi{\linewidth};
    \item mengendalikan kotak, kolom, float, dan posisi elemen secara aman;
    \item mendiagnosis \emph{overfull box} dan memeriksa layout secara visual.
  \end{enumerate}
\end{frame}

\begin{frame}{Peta belajar}
  \tableofcontents
\end{frame}

\section{Model halaman}

\begin{frame}{Layout adalah sistem koordinat, bukan kumpulan spasi}
\small
\begin{columns}[T,onlytextwidth]
\begin{column}{.55\textwidth}
\centering
\begin{tikzpicture}[x=0.72cm,y=0.72cm,>=Latex]
  \fill[codeink!5] (0,0) rectangle (5,7);
  \draw[thick,codeink!70] (0,0) rectangle (5,7);
  \fill[teal!16] (1,1) rectangle (4.2,6.1);
  \draw[dashed,teal!70!black] (1,1) rectangle (4.2,6.1);
  \draw[<->] (0,-.35)--node[below,font=\scriptsize]{margin kiri} (1,-.35);
  \draw[<->] (4.2,-.35)--node[below,font=\scriptsize]{margin kanan} (5,-.35);
  \draw[<->] (0,7.35)--node[above,font=\scriptsize]{lebar kertas} (5,7.35);
  \draw[<->] (4.55,1)--node[right,font=\scriptsize,align=left]{area teks} (4.55,6.1);
  \node[font=\scriptsize,align=center] at (2.6,6.45) {header};
  \node[font=\scriptsize,align=center] at (2.6,.62) {footer};
  \draw[decorate,decoration={brace,amplitude=4pt}] (0,0)--(0,1)
    node[midway,left=5pt,font=\scriptsize,align=right]{margin\ bawah};
  \draw[decorate,decoration={brace,amplitude=4pt}] (0,6.1)--(0,7)
    node[midway,left=5pt,font=\scriptsize,align=right]{margin\ atas};
\end{tikzpicture}
\end{column}
\begin{column}{.41\textwidth}
  \small
  \textbf{Empat lapisan yang perlu dibedakan:}
  \begin{enumerate}
    \item kertas: \dimensi{\paperwidth}, \dimensi{\paperheight};
    \item area teks: \dimensi{\textwidth}, \dimensi{\textheight};
    \item isi lokal: \dimensi{\linewidth}, \dimensi{\columnwidth};
    \item jarak internal: padding, indentasi, dan glue.
  \end{enumerate}
  Hindari mengakali margin dengan menambah \dimensi{\vspace} berulang.
\end{column}
\end{columns}
\end{frame}

\begin{frame}[fragile]{Ukuran penting yang dapat dibaca atau dipakai ulang}
\begin{lstlisting}
\paperwidth   \paperheight   % ukuran kertas
\textwidth    \textheight    % area teks utama
\linewidth                    % lebar baris saat ini
\columnwidth                  % lebar kolom atau minipage
\baselineskip                 % jarak dasar antarbaris
\parindent    \parskip       % indentasi dan jarak paragraf
\headheight   \headsep       % tinggi header dan jaraknya
\footskip                     % jarak baseline terakhir ke footer
\end{lstlisting}
\begin{itemize}
  \item \dimensi{\textwidth} biasanya tetap untuk halaman utama.
  \item \dimensi{\linewidth} mengikuti konteks: kolom, daftar, minipage, atau kotak.
  \item Untuk gambar dan tabel di dalam kolom, gunakan \dimensi{\linewidth}, bukan \dimensi{\textwidth}.
\end{itemize}
\end{frame}

\begin{frame}[fragile]{Margin dengan paket \texttt{geometry}}
\footnotesize
\begin{lstlisting}
\documentclass[a4paper,12pt]{article}
\usepackage[
  a4paper,
  top=2.8cm,bottom=2.5cm,
  left=3cm,right=2cm,
  includehead,includefoot,
  bindingoffset=6mm
]{geometry}
\end{lstlisting}
\begin{itemize}
  \item \dimensi{top}, \dimensi{bottom}, \dimensi{left}, dan \dimensi{right} menetapkan margin.
  \item \dimensi{includehead} dan \dimensi{includefoot} memasukkan header/footer dalam perhitungan area teks.
  \item \dimensi{bindingoffset} menambah ruang untuk jilid; gunakan bersama \dimensi{inner}/\dimensi{outer} pada dokumen dua sisi.
\end{itemize}
\tugas{Uji margin 3 cm dan 25 mm; bandingkan area teksnya.}
\end{frame}

\begin{frame}[fragile]{Margin global dan margin lokal}
\begin{lstlisting}
% Pemeriksaan visual saat pengembangan
\usepackage[showframe]{geometry}

% Mengubah beberapa halaman saja
\newgeometry{top=1.5cm,bottom=1.5cm,left=2cm,right=2cm}
Halaman dengan layout khusus.
\restoregeometry

% Jangan mencampur geometry dengan perubahan manual
% pada \hoffset, \voffset, dan \oddsidemargin
\end{lstlisting}
\begin{block}{Prinsip aman}
  Tetapkan geometri di satu tempat. Bila template sudah memuat \texttt{geometry}, ubah opsinya melalui mekanisme template, bukan dengan menimpa panjang internal secara acak.
\end{block}
\end{frame}

\begin{frame}[fragile]{Header, footer, dan ruang vertikal}
\begin{lstlisting}
\usepackage{fancyhdr}
\pagestyle{fancy}
\fancyhf{}
\fancyhead[L]{Judul pendek}
\fancyhead[R]{\leftmark}
\fancyfoot[C]{\thepage}

% Sediakan ruang yang cukup untuk header
\setlength{\headheight}{14pt}
\end{lstlisting}
\begin{itemize}
  \item \dimensi{\headheight}: tinggi isi header, bukan jarak dari tepi kertas.
  \item \dimensi{\headsep}: jarak antara header dan area teks.
  \item \dimensi{\footskip}: jarak dari baseline baris terakhir ke baseline footer.
  \item Jika header terlalu tinggi, perbaiki \dimensi{\headheight} atau gunakan \texttt{geometry}, bukan \dimensi{\vspace} negatif.
\end{itemize}
\end{frame}

\section{Jarak dan padding}

\begin{frame}{Margin, padding, dan gap}
\begin{tikzpicture}[x=1cm,y=1cm]
  \fill[orange!12] (0,0) rectangle (10,4.2);
  \draw[orange!70!black,thick] (0,0) rectangle (10,4.2);
  \fill[teal!17] (1.25,.8) rectangle (8.75,3.4);
  \draw[teal!70!black,thick] (1.25,.8) rectangle (8.75,3.4);
  \fill[violet!18] (2.1,1.45) rectangle (7.9,2.75);
  \draw[violet!70!black,thick] (2.1,1.45) rectangle (7.9,2.75);
  \node[font=\small] at (5,3.85) {margin: ruang di luar objek};
  \node[font=\small] at (5,1.1) {padding: ruang di dalam kotak};
  \node[font=\small] at (5,2.1) {isi};
  \draw[<->,>=Latex] (8.75,3.65)--node[above,font=\scriptsize]{margin} (10,3.65);
  \draw[<->,>=Latex] (7.9,3.05)--node[above,font=\scriptsize]{padding} (8.75,3.05);
  \draw[<->,>=Latex] (2.1,1.2)--node[below,font=\scriptsize]{padding} (7.9,1.2);
\end{tikzpicture}
\medskip
\begin{itemize}
  \item \textbf{Margin}: jarak objek dari tepi halaman atau objek lain.
  \item \textbf{Padding}: jarak isi dari tepi kotaknya.
  \item \textbf{Gap}: jarak antara dua objek yang berdampingan.
\end{itemize}
\end{frame}

\begin{frame}[fragile]{Padding pada kotak: \texttt{fboxsep} dan \texttt{tcolorbox}}
\begin{lstlisting}
% Kotak LaTeX sederhana
\setlength{\fboxsep}{8pt}       % padding semua sisi
\setlength{\fboxrule}{0.5pt}    % tebal garis
\fbox{\parbox{.72\linewidth}{Isi kotak yang dapat membungkus baris.}}

% Kotak dengan padding per sisi
\begin{tcolorbox}[
  left=8mm,right=8mm,top=3mm,bottom=3mm,
  boxrule=.4pt,arc=1mm]
Isi kotak dengan pengaturan yang lebih terukur.
\end{tcolorbox}
\end{lstlisting}
\begin{itemize}
  \item \dimensi{\fboxsep} berlaku untuk \dimensi{\fbox}, \dimensi{\framebox}, dan beberapa kotak standar.
  \item Untuk kotak berulang, pusatkan aturan pada gaya atau environment agar konsisten.
\end{itemize}
\end{frame}

\begin{frame}[fragile]{Indentasi dan jarak antarparagraf}
\footnotesize
\begin{lstlisting}
% Pola A: paragraf menjorok, tanpa jarak ekstra
\setlength{\parindent}{1.5em}
\setlength{\parskip}{0pt}

% Pola B: tanpa indentasi, ada jarak antarparagraf
\setlength{\parindent}{0pt}
\setlength{\parskip}{6pt plus 2pt minus 1pt}
\end{lstlisting}
\begin{block}{Pilih satu pola}
  Indentasi dan \dimensi{\parskip} sama-sama memberi sinyal awal paragraf; gunakan salah satu secara konsisten.
\end{block}
\begin{alertblock}{Kesalahan umum}
  \texttt{\\} memaksa ganti baris, bukan membuat paragraf. Paragraf baru memakai baris kosong atau \texttt{\\par}.
\end{alertblock}
\end{frame}

\begin{frame}[fragile]{Jarak antarbaris dan glue vertikal}
\begin{lstlisting}
\usepackage{setspace}
\setstretch{1.15}       % seluruh dokumen atau satu grup
{\setstretch{1.30} Teks yang lebih renggang.}

\smallskip              % jarak semantik kecil
\medskip
\bigskip
\vspace{8pt}            % jarak eksplisit
\vspace*{8pt}           % tetap dipertahankan di awal/akhir halaman
\addvspace{8pt}         % tidak menumpuk berlebihan dengan jarak serupa
\end{lstlisting}
\begin{itemize}
  \item Gunakan \dimensi{\smallskip}/\dimensi{\medskip}/\dimensi{\bigskip} untuk hubungan logis.
  \item Gunakan \dimensi{\vspace} untuk koreksi lokal yang benar-benar diperlukan.
  \item Spasi vertikal dapat menghilang di batas halaman; \dimensi{\vspace*} memaksanya tetap ada.
\end{itemize}
\end{frame}

\begin{frame}[fragile]{Jarak pada float dan caption}
\begin{lstlisting}
\setlength{\intextsep}{10pt plus 2pt minus 2pt}
\setlength{\textfloatsep}{12pt plus 2pt minus 2pt}
\setlength{\floatsep}{8pt plus 2pt minus 2pt}
\setlength{\abovecaptionskip}{4pt}
\setlength{\belowcaptionskip}{0pt}
\end{lstlisting}
\begin{itemize}
  \item \dimensi{\intextsep}: jarak float yang berada di tengah teks.
  \item \dimensi{\textfloatsep}: jarak float di atas/bawah halaman dengan teks.
  \item \dimensi{\floatsep}: jarak antara dua float.
  \item Lebih baik mengubah panjang ini daripada menambah \dimensi{\vspace} di setiap gambar.
\end{itemize}
\end{frame}

\begin{frame}[fragile]{Jarak horizontal dan glue fleksibel}
\begin{lstlisting}
Teks\hspace{1em}teks\quad teks\qquad teks

Kiri\hfill Kanan

% Glue dengan lebar alami, dapat meregang dan menyusut
\hspace{1em plus .5em minus .2em}
\hfill
\end{lstlisting}
\begin{itemize}
  \item \dimensi{\hspace} adalah jarak eksplisit; \dimensi{\hfill} adalah ruang yang meregang sampai batas baris.
  \item TeX mengisi ruang dengan \emph{glue}: ukuran alami + kemampuan meregang + kemampuan menyusut.
  \item Karena itu, gunakan panjang elastis untuk distribusi; jangan membuat setiap posisi dengan angka mutlak.
\end{itemize}
\end{frame}

\section{Kotak, kolom, dan struktur lokal}

\begin{frame}[fragile]{\texttt{minipage}: area teks mandiri}
\begin{lstlisting}
\begin{minipage}[t]{.46\linewidth}
  \textbf{Kolom kiri}\par
  Lebar mengikuti konteks saat ini.
\end{minipage}\hfill
\begin{minipage}[t]{.46\linewidth}
  \textbf{Kolom kanan}\par
  Tinggi mengikuti isi kolom.
\end{minipage}
\end{lstlisting}
\begin{itemize}
  \item \dimensi{[t]} menyelaraskan bagian atas; gunakan \dimensi{[c]} atau \dimensi{[b]} bila diperlukan.
  \item Di dalam \texttt{minipage}, \dimensi{\linewidth} berubah menjadi lebar minipage.
  \item \dimensi{\hfill} di antara dua minipage mengisi ruang yang tersedia.
\end{itemize}
\tugas{Buat dua kotak dengan lebar .45\dimensi{\linewidth}; tambahkan tabel kecil hanya pada kolom kanan.}
\end{frame}

\begin{frame}[fragile]{Lebar relatif: \texttt{linewidth} mengalahkan \texttt{textwidth}}
\begin{lstlisting}
\begin{minipage}{.48\textwidth}
  \includegraphics[width=\linewidth]{gambar.pdf}
\end{minipage}

\begin{tabularx}{\linewidth}{@{}>{\raggedright\arraybackslash}Xr@{}}
  Deskripsi fleksibel & Nilai \\
\end{tabularx}
\end{lstlisting}
\begin{block}{Aturan praktis}
  Gunakan \dimensi{\textwidth} untuk ukuran halaman utama; gunakan \dimensi{\linewidth} untuk isi yang harus mengikuti wadah saat ini.
\end{block}
\end{frame}

\begin{frame}[fragile]{Daftar: spacing tanpa mengubah isi}
\begin{lstlisting}
\usepackage{enumitem}
\begin{itemize}[leftmargin=*,itemsep=2pt,topsep=3pt]
  \item Item pertama.
  \item Item kedua dengan jarak yang terukur.
\end{itemize}

% Aturan global yang dapat ditinjau ulang
\setlist[enumerate]{itemsep=2pt,parsep=0pt}
\end{lstlisting}
\begin{itemize}
  \item \dimensi{leftmargin}: posisi label dan isi terhadap batas kiri.
  \item \dimensi{itemsep}: jarak antaritem; \dimensi{topsep}: jarak daftar dengan teks sekitarnya.
  \item Opsi lokal lebih mudah dipelihara daripada banyak \dimensi{\vspace} di antara item.
\end{itemize}
\end{frame}

\begin{frame}[fragile]{Kotak horizontal dan vertikal untuk koreksi kecil}
\begin{lstlisting}
\makebox[.8\linewidth][l]{Judul rata kiri dalam lebar tetap}
\makebox[.8\linewidth][c]{Judul rata tengah}

\raisebox{2pt}[0pt][0pt]{Teks dinaikkan 2 pt}
\phantom{teks}      % ruang seukuran teks, tetapi tidak terlihat
\llap{label}        % menjorok ke kiri tanpa menggeser isi berikutnya
\rlap{label}        % menjorok ke kanan tanpa menggeser isi berikutnya
\end{lstlisting}
\begin{alertblock}{Gunakan dengan hemat}
  Perintah ini berguna untuk label, ikon, dan penyelarasan kecil. Jika dipakai untuk menyusun seluruh halaman, struktur dokumen biasanya perlu diperbaiki.
\end{alertblock}
\end{frame}

\section{Manipulasi layout}

\begin{frame}[fragile]{Kontrol pemutusan halaman}
\begin{lstlisting}
\pagebreak       % minta pemutusan halaman
\newpage         % mulai halaman baru
\clearpage       % kosongkan semua float, lalu halaman baru
\cleardoublepage % halaman kanan pada dokumen twoside

% Hindari pemutusan buruk di sekitar judul
\nopagebreak[4]
\needspace{5\baselineskip} % dari paket needspace
\end{lstlisting}
\begin{itemize}
  \item \dimensi{\newpage} tidak memaksa float keluar; \dimensi{\clearpage} melakukannya.
  \item Jangan menambahkan \dimensi{\newpage} setelah setiap bagian hanya karena satu halaman tampak kurang penuh.
  \item Untuk blok yang harus tetap utuh, sediakan ruang terlebih dahulu atau gunakan environment yang sesuai.
\end{itemize}
\end{frame}

\begin{frame}[fragile]{Float: biarkan LaTeX mencari posisi yang baik}
\begin{lstlisting}
\begin{figure}[htbp]
  \centering
  \includegraphics[width=.75\linewidth]{hasil.pdf}
  \caption{Hasil pengukuran.}
  \label{fig:hasil}
\end{figure}
\end{lstlisting}
\begin{itemize}
  \item \dimensi{h}: here, \dimensi{t}: top, \dimensi{b}: bottom, \dimensi{p}: halaman float.
  \item Urutan \dimensi{[htbp]} adalah preferensi, bukan koordinat mutlak.
  \item Jika gambar harus benar-benar berada di lokasi tertentu, tanyakan dulu apakah float memang tepat untuk elemen tersebut.
\end{itemize}
\end{frame}

\begin{frame}[fragile]{Posisi absolut dengan TikZ: cocok untuk anotasi}
\begin{lstlisting}
\begin{tikzpicture}[remember picture,overlay]
  \node[anchor=north east,xshift=-8mm,yshift=-8mm]
    at (current page.north east) {Catatan};
\end{tikzpicture}
\end{lstlisting}
\begin{itemize}
  \item \dimensi{remember picture} memungkinkan TikZ merujuk koordinat halaman.
  \item \dimensi{overlay} membuat gambar tidak mengambil ruang layout biasa.
  \item Gunakan untuk watermark, nomor, panah, atau anotasi; jangan gunakan untuk paragraf utama yang harus mengalir.
\end{itemize}
\begin{block}{Pertanyaan desain}
  Apakah objek ini bagian dari aliran teks, atau hanya lapisan visual di atasnya? Jawaban itu menentukan apakah kita memakai layout normal atau posisi absolut.
\end{block}
\end{frame}

\begin{frame}[fragile]{Mendiagnosis baris yang keluar margin}
\footnotesize
\begin{lstlisting}
% Tanda hitam pada baris overfull saat debugging
\overfullrule=5pt

% Beri kelonggaran kecil pada pemenggalan paragraf
\emergencystretch=2em

% Pilihan terakhir untuk paragraf yang sulit
\raggedright
\sloppy
\end{lstlisting}
\begin{itemize}
  \item \texttt{Overfull hbox} berarti isi terlalu lebar; cari kata, URL, tabel, atau rumus penyebabnya.
  \item Pendekkan teks, izinkan pemenggalan, gunakan \dimensi{\url}, atau ubah lebar kolom.
  \item \dimensi{\sloppy} adalah pilihan terakhir karena dapat merusak jarak antarkata.
\end{itemize}
\end{frame}

\begin{frame}[fragile]{Pemeriksaan visual layout}
\begin{lstlisting}
\usepackage[showframe]{geometry} % bingkai area halaman
\usepackage{layout}
\begin{document}
\layout                       % diagram ukuran halaman
\end{document}
\end{lstlisting}
\begin{enumerate}
  \item Periksa PDF pada ukuran sebenarnya, bukan hanya pada zoom editor.
  \item Baca log untuk \texttt{Overfull} dan \texttt{Underfull}.
  \item Uji halaman kiri dan kanan untuk dokumen \texttt{twoside}.
  \item Hapus \texttt{showframe}, \texttt{layout}, dan aturan debug sebelum rilis.
\end{enumerate}
\end{frame}

\section{Khusus Beamer}

\begin{frame}{Beamer: area teks juga memiliki margin}
\begin{columns}[T,onlytextwidth]
\begin{column}{.52\textwidth}
\begin{block}{Ukuran global frame}
  \begin{itemize}
    \item margin teks kiri dan kanan;
    \item tinggi frametitle, headline, dan footline;
    \item area isi yang tersisa untuk frame.
  \end{itemize}
\end{block}
\end{column}
\begin{column}{.43\textwidth}
  Jangan memakai \texttt{geometry} untuk mengubah margin Beamer secara sembarang. Beamer memiliki mekanisme sendiri untuk isi frame dan template presentasi.
\end{column}
\end{columns}
\end{frame}

\begin{frame}[fragile]{Mengatur margin dan alignment pada Beamer}
\begin{lstlisting}
\setbeamersize{text margin left=8mm,text margin right=8mm}
\end{lstlisting}
Untuk frame tertentu, tambahkan opsi \texttt{[t]} agar isi dimulai dari atas.
Frame berisi kode verbatim menggunakan opsi \texttt{[fragile]}.
\begin{itemize}
  \item \dimensi{[t]}, \dimensi{[c]}, dan \dimensi{[b]} mengatur alignment vertikal isi frame.
  \item Frame berisi \texttt{verbatim} atau \texttt{lstlisting} membutuhkan \dimensi{[fragile]}.
  \item Uji margin pada frame judul, frame dua kolom, dan frame kode; hasilnya tidak selalu sama.
\end{itemize}
\end{frame}

\begin{frame}[fragile]{Kolom Beamer: gunakan \texttt{onlytextwidth}}
\begin{lstlisting}
\begin{columns}[T,onlytextwidth]
  \begin{column}{.48\textwidth}
    Isi kiri; \linewidth menyesuaikan kolom.
  \end{column}
  \begin{column}{.48\textwidth}
    \includegraphics[width=\linewidth]{hasil.pdf}
  \end{column}
\end{columns}
\end{lstlisting}
\begin{itemize}
  \item \dimensi{[T]} menyelaraskan bagian atas kolom.
  \item \dimensi{onlytextwidth} membuat kolom memakai lebar area teks, bukan menambah ruang ekstra.
  \item Dua kolom .48\dimensi{\textwidth} menyisakan ruang untuk gap di tengah.
\end{itemize}
\end{frame}

\begin{frame}[fragile]{Ruang elastis di dalam frame}
\begin{lstlisting}
  \textbf{Pesan utama}
  \medskip
  Penjelasan singkat.
  \vfill
  \begin{alertblock}{Kesimpulan}
    Isi ini didorong ke bagian bawah frame.
  \end{alertblock}
\end{lstlisting}
\begin{itemize}
  \item \dimensi{\vfill} mengisi sisa ruang vertikal dan berguna untuk memisahkan pesan utama dari kesimpulan.
  \item Jika slide terlalu penuh, kurangi isi atau pecah menjadi dua frame; jangan hanya mengecilkan font.
  \item Jarak kosong adalah alat hierarki: ia mengelompokkan isi yang berhubungan.
\end{itemize}
\end{frame}

\section{Latihan dan ringkasan}

\begin{frame}{Latihan inti: bedah layout}
  Ambil satu dokumen artikel atau laporan yang pernah dibuat.
  \begin{enumerate}
    \item Tampilkan frame halaman dengan \texttt{showframe}; catat ukuran kertas, margin, dan area teks.
    \item Buat satu kotak dengan padding berbeda untuk atas/bawah dan kiri/kanan.
    \item Susun dua minipage dengan lebar relatif dan gambar selebar \texttt{\\linewidth}.
    \item Atur jarak paragraf atau jarak antaritem tanpa menambah \texttt{\\vspace} berulang.
    \item Sengaja buat satu \texttt{Overfull hbox}; diagnosis dan perbaiki penyebabnya.
  \end{enumerate}
  \tugas{Bandingkan hasil sebelum dan sesudah. Jelaskan perubahan dengan nama panjang LaTeX yang tepat.}
\end{frame}

\begin{frame}{Ringkasan: urutan memilih alat}
\footnotesize
\begin{enumerate}
  \item Atur kertas dan margin dengan \texttt{geometry}.
  \item Atur header, footer, dan page style dengan paket yang sesuai.
  \item Atur paragraf, daftar, dan float melalui panjang atau opsi environment.
  \item Gunakan \texttt{minipage}, \texttt{tabularx}, dan \texttt{tcolorbox} untuk struktur lokal.
  \item Pilih \texttt{\\linewidth} agar isi mengikuti wadah.
  \item Gunakan \texttt{\\vspace}, \texttt{\\hspace}, dan posisi absolut hanya sebagai koreksi lokal.
  \item Kompilasi, baca log, dan periksa PDF pada ukuran sebenarnya.
\end{enumerate}
\begin{block}{Prinsip akhir}
  Layout kuat berasal dari struktur dan ukuran relatif; angka manual dipakai secara lokal.
\end{block}
\end{frame}

\begin{frame}[plain]
  \centering
  {\usebeamerfont{title}Pertanyaan dan diskusi\par}
  \medskip
  Margin mana yang ingin Anda kendalikan lebih dulu?
\end{frame}

\end{document}
